\section{Local Half-Split Certificate}
\label{sec:unified-reality-local-half-split-certificate}

The resonance picture should not be read as if an entire infinite prime
spiral suddenly collapses to zero.  A zero-relevant event is local: after
dominance and normalization have been displayed, the completed read is
cut by a finite two-branch interference island.  The RH-facing content is
that a stable local cut must be a half-split.  The real coordinate
controls length balance; the imaginary coordinate records the remaining
phase quantization.

\begin{theorem}[Two-vector silence is equal length and opposite phase]
\label{thm:unified-reality-two-vector-silence-equal-length-opposite-phase}
For nonzero complex vectors $A,B\in\ComplexUp$, the row
$$
A+B=0
$$
is equivalent to
$$
|A|=|B|
$$
and
$$
\arg(A)-\arg(B)\equiv\pi\pmod{2\pi}.
$$
\end{theorem}

\begin{proof}
If $A+B=0$, then $A=-B$, so the two vectors have equal modulus and
opposite phase.  Conversely, equal modulus and phase difference $\pi$
give $A=-B$, hence $A+B=0$.
\end{proof}

\begin{definition}[Local two-branch zero island]
\label{def:unified-reality-local-two-branch-zero-island}
A \emph{local two-branch zero island} for a critical-strip row $s$ is a
local decomposition
$$
\Xi(u)=A_n(u)+B_n(u)+R_n(u)
$$
on a disk $D_n$ around $s$, together with rows saying that $A_n$ and
$B_n$ are the two dominant branches and that the remainder satisfies the
Rouch\'e boundary discipline required to preserve the local zero count.
The row is zero-relevant only after these dominance and tail rows are
displayed; a visible two-term split without them is merely a projection
model.
\end{definition}

\begin{definition}[Half-split coordinate]
\label{def:unified-reality-half-split-coordinate}
For a local two-branch zero island, the \emph{half-split coordinate} at
$s$ is the length-balance row
$$
|A_n(s)|=|B_n(s)|.
$$
Equivalently, the zero cuts the normalized local segment into two equal
halves.  The complementary phase row is
$$
\arg(A_n(s)/B_n(s))\equiv(2k_n-1)\pi
\pmod{2\pi}.
$$
The half-split row records the real-length constraint, while the phase
row records the height or time quantization.
\end{definition}

\begin{definition}[Real-defect split formula]
\label{def:unified-reality-real-defect-split-formula}
A local two-branch zero island has a \emph{real-defect split formula}
when there is a positive branch-separation scale $H_n$ and a bounded
correction row $\varepsilon_n$ such that
$$
\log\frac{|A_n(s)|}{|B_n(s)|}
=
-2\delta H_n+\varepsilon_n
$$
for
$$
s=\frac12+\delta+i\tau.
$$
The term $-2\delta H_n$ is the normal length bias.  The correction
$\varepsilon_n$ belongs to amplitude normalization, finite-window tail
control, or coordinate choice; it is not allowed to encode a hidden
noncompact real scale along a cofinal tower.
\end{definition}

\begin{definition}[Zeta local half-split certificate]
\label{def:unified-reality-zeta-local-half-split-certificate}
For a critical-strip zeta zero $s$, a \emph{zeta local half-split
certificate} consists of a cofinal family of local two-branch zero
islands with:
$$
\mathsf{DominantTwoBranch}(A_n,B_n),
\qquad
\mathsf{RoucheTailControl}(R_n),
$$
$$
\mathsf{HalfSplit}(A_n,B_n,s),
\qquad
\mathsf{RealDefectSplitFormula}(H_n,\varepsilon_n),
$$
and a no-hidden-real-correction row saying that $\varepsilon_n$ cannot
cofinally cancel $-2\delta H_n$ when $\delta\#0$ and the displayed
branch-separation scales remain unbounded or otherwise non-negligible.
\end{definition}

\begin{theorem}[Local half-split certificates are conditionally sufficient]
\label{thm:unified-reality-local-half-split-certificates-conditionally-sufficient}
If every critical-strip zeta zero carries a zeta local half-split
certificate, then the RH fixed-half section exists.
\end{theorem}

\begin{proof}
Let $z:\mathcal Z_{\zeta}$ have complex coordinate
$$
s=\frac12+\delta+i\tau.
$$
The certificate supplies local islands with
$$
|A_n(s)|=|B_n(s)|.
$$
Thus
$$
\log\frac{|A_n(s)|}{|B_n(s)|}=0.
$$
The real-defect split formula gives
$$
0=-2\delta H_n+\varepsilon_n.
$$
If $\delta\#0$, the normal term cannot be cofinally cancelled by
$\varepsilon_n$ by the no-hidden-real-correction row.  This contradicts
the half-split equation.  Hence $\neg(\delta\#0)$, and located
fixed-half collapse gives $\delta=0$.  Therefore $\Re(s)=1/2$, and the
fixed-half section follows from
\autoref{thm:unified-reality-rh-fixed-half-section-equivalence}.
\end{proof}

\begin{corollary}[Phase carries the remaining zero data]
\label{cor:unified-reality-phase-carries-remaining-zero-data}
Once a zero-relevant local island has been normalized to a half-split,
the remaining zero-location data are phase data.
\end{corollary}

\begin{proof}
By
\autoref{thm:unified-reality-two-vector-silence-equal-length-opposite-phase},
the local zero condition for the dominant pair is equal length together
with opposite phase.  The half-split certificate supplies the equal
length row.  The remaining equation is therefore
$$
\arg(A_n(s)/B_n(s))\equiv(2k_n-1)\pi
\pmod{2\pi}.
$$
If the active branches have phase factors of the form
$$
A_n(s)\sim C_+e^{-i\tau H_+},
\qquad
B_n(s)\sim C_-e^{-i\tau H_-},
$$
then the height coordinate $\tau$ is determined by the corresponding
phase congruence.  The real part supplies length balance; the imaginary
part records phase quantization.
\end{proof}

\begin{corollary}[Quartic endpoint zeros supply the finite template]
\label{cor:unified-reality-quartic-endpoint-half-split-template}
The fixed quartic endpoint-zero law is the finite model for the local
half-split certificate.
\end{corollary}

\begin{proof}
In the quartic carrier, endpoint zeros are obtained only after the
positive-axis dominant branches, local Puiseux expansions,
partial-fraction amplitudes, and Rouch\'e tail control are displayed, as
recorded in
\autoref{thm:unified-reality-quartic-endpoint-zero-branch-interference}.
The active two-branch sum is transported to a shifted cosine model, whose
zeros occur at odd half-periods.  This is precisely the finite
half-split pattern: equal dominant lengths with phase difference an odd
multiple of $\pi$.  The quartic result remains a fixed-window finite
template, not a global zeta theorem.
\end{proof}

\begin{corollary}[Local half-split refines resonance no-hair]
\label{cor:unified-reality-local-half-split-refines-resonance-no-hair}
The local half-split certificate refines the resonance no-hair boundary
by identifying the finite zero island on which the normal-energy
condition is zero-relevant.
\end{corollary}

\begin{proof}
The resonance no-hair boundary of
\autoref{def:unified-reality-resonance-no-hair-boundary} states that a
stable zero-relevant resonance family has zero normalized detuning.  The
local half-split certificate explains what makes the family
zero-relevant: a two-branch dominant island, Rouch\'e tail control,
length equality, and a real-defect split formula.  Thus the certificate
prevents projected silence from being mistaken for a proof object; it
requires the local segment on which the zero actually acts.
\end{proof}

\begin{corollary}[Relation to solenoid and selector routes]
\label{cor:unified-reality-local-half-split-relation-to-solenoid-selector-routes}
The split-prime, Dedekind, and Zeckendorf solenoid routes supply possible
source mechanisms for the branch-separation scale $H_n$ in a local
half-split certificate.
\end{corollary}

\begin{proof}
In the fixed split-prime route, $H_n$ is supplied by the repeated local
rows above a rational prime $q$.  In the Dedekind route, it is supplied
by the prime-ideal rows in the lifted tower.  In the Zeckendorf selector
route, it is supplied by the golden-cofinal prime family selected by the
no-adjacent grammar.  Each route contributes a source for the length
ledger; solenoid compatibility and Lee--Yang stability are then the
candidate mechanisms preventing a hidden real correction from absorbing
a nonzero $\delta$.
\end{proof}
